\section{Introduction}
\label{sec:introduction}

Voice interfaces in vehicles must operate in an acoustic environment that differs substantially from the clean, close-talk speech commonly used to develop speech models. Road and powertrain noise, ventilation, reverberation, loudspeaker playback, microphone placement, and speech from other occupants can all affect the signal captured by an in-cabin microphone. Recent in-car corpora such as AISHELL-5 and CAVEMOVE make these conditions more accessible for systematic research: AISHELL-5 provides real multichannel in-car speech and noise, while CAVEMOVE provides vehicle impulse responses and noise recordings that support the synthesis of cabin mixtures \cite{aishell5,cavemove}. These resources also highlight that performance on generic speech mixtures alone does not establish robustness inside a vehicle.

A voice assistant must solve more than acoustic separation. Given a target speaker, the system should retain that speaker's speech while suppressing competing talkers and background sound, determine when the target is speaking, detect an activation phrase, interpret the following command, and map the inferred intent to an appropriate action. These stages address distinct questions: who is speaking, whether the intended user is speaking now, whether the system should activate, and what operation is requested. A conventional voice activity detector answers whether speech is present, but does not by itself establish that the speech belongs to the enrolled user. Target-speaker activity modeling has therefore been studied as a speaker-conditioned alternative for determining speaker activity over time \cite{tsvad}. For an in-car assistant, this distinction matters because a passenger speaking a wake phrase should not necessarily activate a driver-oriented command path.

Target speaker extraction (TSE) provides a useful front end for this setting. Unlike speaker-independent separation, which estimates a set of sources without knowing which output corresponds to the desired person, TSE uses a reference utterance to condition extraction on the target identity. SpEx+ is a time-domain TSE architecture that uses a shared representation for mixture and reference speech and achieved strong results on standard speech-mixture benchmarks \cite{spexplus}. However, success on generic or synthetically mixed speech does not guarantee performance under vehicle cabin acoustics. Domain mismatch between training mixtures and in-car recordings motivates evaluating and adapting extraction with cabin-specific acoustic conditions. LibriSpeech and Libri2Mix offer scalable clean speech and speech-mixture resources for model development, while CAVEMOVE and vehicle-interior recordings can contribute driving noise and cabin acoustics \cite{librispeech,librimix,cavemove,visdataset}. AISHELL-5 further provides real multichannel, multi-speaker in-car recordings for speech technology research \cite{aishell5}.

The downstream stages also require task-specific data and evaluation. Keyword spotting detects a short, predefined phrase in an audio stream; Google Speech Commands was introduced as a benchmark for limited-vocabulary recognition and keyword spotting \cite{speechcommands}. Spoken language understanding maps an utterance to a structured intent, and the Fluent Speech Commands corpus supports evaluation of direct speech-to-intent systems \cite{fluentcommands}. Yet high accuracy on an isolated, clean intent test set does not alone measure whether the intended in-car user was selected, whether an activation phrase was accepted correctly, or whether the complete pipeline executed the right action. It is therefore useful to describe and assess the system both by stage and as a cascade, with the limits of each evaluation made explicit.

This work presents a multi-stage in-car voice-processing pipeline built around a SpEx+-based target speaker extractor adapted using vehicle-cabin noise. The extracted signal is passed through a target-speech presence gate, wake-word detection, intent classification, and intent-to-action mapping. The stages use datasets suited to their respective learning problems: LibriSpeech and Libri2Mix for speech and mixture training; CAVEMOVE, AISHELL-5, and the Vehicle Interior Sound Dataset for cabin acoustics, noise, and evaluation or augmentation; Google Speech Commands V2 for keyword spotting; and Fluent Speech Commands for intent classification. The pipeline is designed to make the activation and command path conditional on the target speaker, reducing the risk that another occupant's speech triggers an action.

On the reported held-out CAVEMOVE evaluation, the adapted extraction stage improves mean SI-SDR by 10.98 dB relative to the mixture input. The intent classifier achieves 99.31\% top-1 accuracy on the clean Fluent Speech Commands test set. These figures demonstrate promising performance for two individual stages under their stated evaluation conditions; they should not be interpreted as an end-to-end wake-to-action accuracy. The contribution of this work is the construction and stage-wise evaluation of a cabin-adapted, target-conditioned voice pipeline that connects extraction to target-speech gating, wake-word detection, intent recognition, and action mapping. The remainder of this paper describes the datasets and methods, reports stage-specific and pipeline evaluations, and discusses the limitations and requirements for validating the complete system under real in-car use.
